\section{Quantum Phase Estimation}
One of the most important quantum algorithms to date is the quantum phase estimation \texttt{QPE}.
It allows to measure the eigenvalues of unitary matrices
\begin{align}
    U \ket{\alpha} = e^{i\alpha} \ket{\alpha}, \;\; \alpha \in \mathbb{R}
\end{align}
and by extension also the eigenvalues of Hermitian matrices
\begin{align}
       H \ket{\alpha} = \alpha \ket{\alpha}, \;\; \alpha \in \mathbb{R} 
\end{align}
through the exponential connection $U=e^{iH}$.~\footnote{In a way the \texttt{QPE} gives an explicit implementation of the quantum mechanical axiom that observables are described through Hermitian matrices, with their eigenvalues being the observable quantities.}

As we want to use qubits to ``read-out'' our eigenvalues $\alpha$ we need to make some approximations as real numbers can not be exactly described by a finite sequence of bits.
For now we will operate under assumption that $\alpha = 2\pi \frac{k}{2^n}$ with integers $n \in \mathbb{N}$ and $0 \leq k < 2^n$, meaning that we assume our eigenvalues can be exactly represented with $n$ bits where $k$ is an integer represented by the bitstring $\mathbf{k}$.

\begin{algo*}[\texttt{QPE}: High-Level Description]\phantom{...}\\
Given a Hermitian $H$ with eigenstates $\ket{\alpha}$, eigenvalues $\alpha = 2\pi \frac{k}{2^n}$, and $k$ being a integer represented by bitstring $\mathbf{k}$, then
\texttt{QPE} will act on eigenstates $\ket{\alpha}$ as
\begin{align}
    \texttt{QPE} \left(\ket{\mathbf{0}} \otimes \ket{\alpha}\right) = \ket{\mathbf{k}} \otimes \ket{\alpha}.
\end{align}
\end{algo*}
The \texttt{QPE} operates on two qubit registers: one representing the state of interest $\ket{\alpha}$ (the state register) and another auxilliary $n$-qubits (the read-out register).

Measurement of the first qubit register now gives back the information about the phase $\alpha = 2\pi \frac{k}{2^n}$ by providing $k$ in binary. Note that the eigenstate in the second register remains unchanged -- in fact, the form $\ket{\psi} \otimes \ket{\alpha}$, with $\ket{\psi}$ being some state on the first register, remains intact throughout the whole algorithm.

In the following we are gonna approach the \texttt{QPE} in a step-by-step approach figuring out how to construct it.

\subsection{Single Qubit Case}
In the case that $n=1$, the \texttt{QPE} can be implemented as 
\begin{align}
    \texttt{QPE} = H(q_1) CU(q_1,q_2) H(q_1) 
\end{align}
or as circuit diagram (also indicating where to read-out after the unitaries are executed):
\begin{center}
$ \texttt{QPE}$ = \begin{quantikz}
& \gate{H} & \ctrl{1}& \gate{H} & \meter{} & \\
& \hphantomgate{} & \gate{U} &  \hphantomgate{} &\qw  
\end{quantikz}
\end{center}
Let's see what this circuit computes when applied to an eigenstate of $U$
\begin{align}
    \texttt{QPE}\left(\ket{0}\otimes\ket{\alpha}\right) &= \frac{1}{\sqrt{2}} H(q_1) \left( CU(q_1,q_2) \ket{0}\otimes\ket{\alpha} + CU(q_1,q_2)\ket{1}\otimes\ket{\alpha}\right) \\
    &= \frac{1}{\sqrt{2}} H(q_1) \left( \ket{0}\otimes\ket{\alpha} + e^{i\alpha}\ket{1}\otimes\ket{\alpha}\right) \\
    &= \frac{1}{\sqrt{2}} H(q_1) \left( \ket{0} + e^{2\pi i \frac{k}{2} }\ket{1}\right)\otimes\ket{\alpha}\\
    &= \frac{1}{\sqrt{2}} \left( \ket{+} + e^{2\pi i \frac{k}{2} }\ket{-}\right)\otimes\ket{\alpha}\label{eq:qpe-1qubit}
\end{align}
now remember that by our initial assumption, we are only considering matrices $U$ with eigenvalues exactly representable as $\alpha = 2\pi k/2^n$ -- so for $n=1$ we have two possibilities $k=0$ and $k=1$
\begin{align}
    \texttt{QPE}\left(\ket{0}\otimes\ket{\alpha}\right) &= \frac{1}{\sqrt{2}} 
    \begin{cases}
    \ket{+} + \ket{-} \otimes \ket{\alpha} \text{ for } k=0 \\
    \ket{+} - \ket{-} \otimes \ket{\alpha} \text{ for } k=1
    \end{cases}
\end{align}
Since we have $1/\sqrt{2}(\ket{+}+\ket{-}) = \ket{0}$ and $1/\sqrt{2}(\ket{+}-\ket{-}) = \ket{1}$, we can summarize as
\begin{align}
   \texttt{QPE}\left(\ket{0}\otimes\ket{\alpha}\right) & = \ket{k}\otimes\ket{\alpha}
\end{align}
so this is exactly what \texttt{QPE} promised to do.

\subsection{2-Qubit Case}
In the two-qubit case we are assuming that our phase takes the discrete values 
\begin{gather}
    e^{i\alpha} = e^{i2\pi k/4} = \left\{ 1, i, -1, -i \right\} 
\end{gather}
for $k=\left\{0,1,2,3\right\}$.
We are now having the following qubit registers: register $q_1$ consisting of $n=2$ qubits which we will just label as $0$ and $1$ and the second register $q_2$ holding the qubits necessary to represent the state $\ket{\alpha}$. The implementation in this case looks as
\begin{align}
    \texttt{QPE} = F(q_1) CU^2(0,q_2) CU(1,q_2) \textsc{H}(0)\textsc{H}(1) 
\end{align}

where the last operation $F(q_1)$ acts on the first qubit register and is at this point undetermined. 
As a circuit diagram that looks as
\begin{center}
$ \texttt{QPE}$ = \begin{quantikz}
& \gate{\hd} & \qw      & \ctrl{2}&  \gate[1,2]{F}  & \meter{} &  \\
& \gate{\hd} & \ctrl{1} & \qw     &  \qw  & \meter{} & \\
& {/}\qw     & \gate{U} & \gate{U^2}&  \qw  &\qw       &    
\end{quantikz}
\end{center}

Let's see how the state looks before $F$ acts. We will label the two qubits in the read-out register as $1$ and $2$ and use the label $q$ for the collection of qubits in the state register
\begin{align}
    \texttt{QPE}\left(\ket{0}\otimes\ket
    {0}\otimes\ket{\alpha}\right) &= F(1,2) CU^2(1,q) CU(1,q) \left(\ket{+}\otimes\ket{+} \otimes \ket{\alpha}\right)\nonumber \\
    & = \frac{1}{\sqrt{2}}F(1,2) CU^2(0,q_2) \ket{+}\otimes\left( \ket{0} + e^{2\pi i \frac{k}{2} }\ket{1}\right) \otimes \ket{\alpha} \nonumber\\
    & = \frac{1}{2}F(1,2) \left( \ket{0} + \left(e^{2\pi i \frac{k}{4} }\right)^2\ket{1}\right) \otimes\left( \ket{0} + e^{2\pi i \frac{k}{4} }\ket{1}\right) \otimes \ket{\alpha} \nonumber\\
    & = \frac{1}{2} F(1,2) \left( \ket{00} +  e^{2\pi i \frac{k}{4} }\ket{01} + e^{2\pi i \frac{2k}{4} }\ket{10} + e^{2\pi i \frac{3k}{4} }\ket{11} \right) \otimes \ket{\alpha}\nonumber \\
    & \text{let's switch to binary notation on the first register}\nonumber\\
    & = \frac{1}{2} F(1,2) \left( e^{2\pi i \frac{\mathbf{0}k}{4}}\ket{\mathbf{0}} +  e^{2\pi i \frac{\mathbf{1}k}{4} }\ket{\mathbf{1}} + e^{2\pi i \frac{\mathbf{2}k}{4} }\ket{\mathbf{2}} + e^{2\pi i \frac{\mathbf{3}k}{4} }\ket{\mathbf{3}} \right) \otimes \ket{\alpha} \nonumber\\
    & = \frac{1}{2} F(0,1) \sum_{l<4}\left( e^{2\pi i \frac{lk}{4}} \ket{l} \right)\otimes\ket{\alpha}.
\end{align}
All that is left now is to determine the action of $F$ on the qubits of the read-out register. We want it to act such that
\begin{align}
     F(0,1)  \frac{1}{2}\sum_{l<4}\left( e^{2\pi i \frac{lk}{4}} \ket{l}\right) = \ket{k}
\end{align}
which is exactly the definition of the discrete Fourier transformation on dimension $N=2^2=4$. In particular here, $F$ should act as the inverse Fourier transformation. If we can define the Fourier transformation \texttt{QFT} on qubits as a unitary operation then $F = \texttt{QFT}^\dagger$. At this point we assume that we have access to \texttt{QFT} and we will see later how to construct it.

\subsection{$N$-Qubit Case}
Generalizing form the 2-qubit case to the $N$-qubit case is now relatively straightforward. We assume now that our phase can be captured by any of the $2^n$ unit roots $e^{2\pi i \frac{k}{2^n}}$ and implement the \texttt{QPE} as follows.
\begin{algo*}[\texttt{QPE} (concrete construction)]
\begin{align}
\label{eq:general_qpe}
    \texttt{QPE} = \textsc{QFT}^\dagger(q_1) \left(\prod_{q<n} CU^{2^q}(q,q_2) \textsc{H}(q)\right),
\end{align}
or as circuit diagram
\begin{center}
$ \texttt{QPE}$ = \begin{quantikz}
& \gate{\hd} & \qw      & \qw       &\qw &{\cdots}& & \ctrl{4}            & \gate[4]{\texttt{QFT}^\dagger}  & \meter{} &  \\
& \vdots & & & & & & & & & \\
& \gate{\hd} & \qw      & \ctrl{2}  &\qw &{\cdots}& & \qw                 &  \qw                                & \meter{} &  \\
& \gate{\hd} & \ctrl{1} & \qw       &\qw &{\cdots}& &\qw                 &  \qw                                & \meter{} & \\
& {/}\qw     & \gate{U} & \gate{U^2}&\qw &{\cdots}& &  \gate{U^{2^{n-1}}}      &  \qw                                &\qw       &    
\end{quantikz}
\end{center}
\end{algo*}
The operation $CU^{2^q}(q,q_2)$ may look expensive, remember however that we can represent unitaries as
\begin{align}
    U(\theta) = e^{-i\frac{\theta}{2} G}
\end{align}
with a Hermitian generator $G$. With this we have
\begin{align}
    U(\theta)^n = U(n\theta).
\end{align}
So $CU^{2^q}(q,q_2)$ does not necessarily mean $2^q$ repetitions of the same primitive operation (however, arbitrary small angles might also come with some computational overhead -- we will however not care about those specifics for now). 

\subsection{Imperfections}
In general, we will not be able to represent an arbitrary real number $0 \leq \alpha < 2\pi$ with a finite precision $n$ as $\alpha = 2\pi \frac{k}{2^n}$ but have to account for a small error $\epsilon \in \mathbb{R}$ for convenience also represented in units of $2\pi$
\begin{align}
    \alpha = 2\pi \frac{k}{2^n} + 2\pi\epsilon.
\end{align}
Note that we can assume $0 \leq |\epsilon| \leq \frac{1}{2^{n+1}}$. Take for example $\epsilon=\frac{1}{2^{n+1}}$, which produces an $\alpha$ whose value exactly between the ones generated with $k$ and $k+1$, so every $\epsilon$ that is larger than $\frac{1}{2^{n+1}}$ should rather be introduced as error for the $\alpha$ value generated by $k+1$.

With this imperfection in our representation of $\alpha$ we can revisit our derivations from before and estimate~\footnote{we have used $e^{2\pi i \epsilon} \approx 1 + 2\pi i \epsilon$ for small $\epsilon$}
\begin{align}
    \texttt{QPE}\ket{0}\otimes\ket{\alpha}
    &\approx \ket{k}\otimes\ket{\alpha} + 2\pi \epsilon i \ket{\phi} \otimes \ket{\alpha}
\end{align}
with $\ket{\phi}$ being and undetermined state.
From this we can see that the failure probability is proportional to the square of our error in the binary representation of the phase. In other words, a finite error $\epsilon$ in the binary representation of $\alpha$ leads to a failure probability of the \texttt{QPE} proportional to $\epsilon^2$. Or with respect to the bit-precision $n$ we have a failure probability proportional to $\epsilon^2 \leq \left(\frac{1}{2^{n+1}}\right)^2$.

Let's go back to the single qubit case. Instead of Eq.~\eqref{eq:qpe-1qubit} we will have
\begin{align}
    \texttt{QPE}\ket{0}\otimes\ket{\alpha} &= \frac{1}{\sqrt{2}}\left(\ket{+} + e^{2\pi\left(\frac{k}{2} + \epsilon\right)}\ket{-}\right) \otimes\ket{\alpha}
\end{align}
The exact failure probability for the case $k=0$ is then
\begin{align}
   \bra{1\alpha}\texttt{QPE}\ket{0\alpha} &=\frac{1}{2} \lvert\braket{1}{+} + e^{2\pi i \epsilon}\braket{1}{-}\rvert^2 \\
   & = \frac{1}{4} \lvert 1 - e^{2\pi i \epsilon} \rvert^2 \\
   & = \frac{1}{2}\left(1-\cos\left(2\pi\epsilon\right)\right).
\end{align}
As expected, there is no possibility of failure in the algorithm for $\epsilon=0$ while the highest probability of failure is $50\%$ for $\epsilon=\frac{1}{2^{1+1}}=\frac{1}{4}$.
Note also that $1-cos(x) \approx 1 - 1 + x^2 = x^2$, so this exact failure probability is consistent with our initial estimate of $\epsilon^2$.
